% FILE: 06_contribution_to_thesis.tex
% Project: atlas
% Role: doctrine-atlas-and-thesis-chapter-surface-map

\section{Contribution to the Dissertation}
\label{sec:atlas-contribution}

\subsection{Contribution to Prediction vs.~Coordination}
\label{sec:atlas-contribution-prediction}

The atlas makes a structural contribution to the thesis argument that process
intelligence is a coordination technology, not a prediction technology. Each of
the seven doctrine modules encodes a coordination law --- a rule about what must
happen between actors, artifacts, and systems at a named boundary --- rather than
a predictive claim about future behavior.

The Aalst Broadcaster exemplifies this. It does not predict the next commit; it
coordinates the transformation of Git history into an OCEL stream so that every
past commit becomes admissible process evidence. The Conformance Stream does not
predict conformance; it continuously coordinates the actual event sequence
against the declared WF-net or Petri Net model and emits severity-rated alerts
when the coordination breaks down. The Blue River Dam operator chain
$\kappa(\rho(\alpha(\mu(O^*))))$ does not predict whether a process instance
will be ALIVE; it enforces the coordination law that every transition must emit
a receipt and every receipt must be gated before the next transition is
authorized.

The AGI Adversarial Defenses module makes this distinction sharpest: the four
Iron Law safeguards (Schr\"odinger's Interface Fallback, Turing Trap Circuit
Breakers, Causal Window P2P Constraints, Hardware Attestation via ZKP) are not
predictive filters; they are coordination boundaries that prevent complexity
collapse regardless of what the AI component predicts. A Turing Trap Circuit
Breaker does not predict infinite reasoning loops; it enforces a boundary that
terminates them.

\textbf{[AUTHOR THESIS]} The atlas thus contributes the claim that the thesis
argument ``prediction does not conserve consequence; coordination does'' is
instantiated across all seven doctrine modules --- each module names a
coordination law that would be dissolved if prediction were substituted for it.

\subsection{Contribution to Process Evidence}
\label{sec:atlas-contribution-evidence}

The atlas contributes to the process evidence strand of the thesis by naming the
\texttt{Evidence<T, State, Witness>} typed receipt structure as the canonical
form of process evidence across all seven modules. This is not a contribution of
the atlas alone --- the SPR thesis formal objects define it, and the
wasm4pm-compat type foundry enforces it at the Rust type level. But the atlas
makes the contribution visible at the doctrine-surface level by showing that
each module produces or consumes typed evidence rather than untyped outputs.

The Aalst Broadcaster produces OCEL-typed event objects. The GGEN Prompt
Manufactory produces Prompt-typed research warrants. The Supabase Law Substrate
produces RLS-typed authorization decisions. The WASM Boundary Law enforces that
the TypeScript projection (Specta) and the binary execution (tsify/wasm-bindgen)
are separately typed surfaces that cannot be silently coerced into each other.
The ZoeApp Ministry Ontology produces process-law-typed lifecycle states
(healthy, degraded, critical). In every case, the type is the evidence: an
untyped output is not admissible as process evidence under the covenant
definition.

The RECEIPT\_REGISTRY defines seven canonical receipts that correspond to the
seven major evidence surfaces of the foundry (papers, PM4Py oracle, gaps,
lifecycle, M\&A, standards, adversarial). The atlas maps its seven doctrine
modules to this receipt structure at the naming level, establishing which
doctrine module relies on which receipt category for its evidentiary basis.

\subsection{Contribution to Manufacturing Architecture}
\label{sec:atlas-contribution-manufacturing}

The atlas contributes to the manufacturing architecture of the dissertation by
providing the naming layer that connects the Chatman Equation
$A = \mu(O^*, T, L)$ to the implementation repositories that instantiate it.
Without the atlas, the manufacturing architecture would require readers to
navigate seven distinct repositories to understand how the four operators
($\mu$, $\alpha$, $\rho$, $\kappa$) are instantiated in each domain. With the
atlas, the mapping is explicit:

\begin{itemize}
  \item $\mu$ is instantiated by the GGEN Prompt Manufactory (TTL $\to$ RQ
        $\to$ TERA $\to$ Prompt) and by the Aalst Broadcaster (Git $\to$ OCEL).
  \item $\alpha$ is instantiated by the Blue River Dam actuation boundary:
        Future Claim (Intent) $\to$ Public Law (Evaluation) $\to$ Actuation
        Boundary (Execution) $\to$ Receipt (Proof).
  \item $\rho$ is instantiated by the WASM Boundary Law (every execution across
        the membrane emits a typed receipt) and by the Supabase RLS Authorization
        Court (every SQL command decision is audited and logged).
  \item $\kappa$ is instantiated by the Blue River Dam gate
        ($\kappa(\tau) \in \{\textsc{ALIVE}, \textsc{PARTIAL}, \textsc{REFUSED}\}$)
        and by the AGI Adversarial Defenses (circuit breakers and fallbacks enforce
        $\kappa$ against adversarial inputs).
\end{itemize}

The atlas also contributes the Combinatorial Maximalism thesis point: the
manufacturing surface is not additive but multiplicative. Adding one standard
(one column) adds $|L| \times |A| \times |F| \times |R| \times |Q| \times |B|
\times |D|$ new manufacturing cells --- not one.

\subsection{Contribution to Capital-Flow Infrastructure}
\label{sec:atlas-contribution-capital}

The atlas contributes to the capital-flow infrastructure strand of the thesis
through two modules: the Supabase Law Substrate and the ZoeApp Ministry
Ontology.

The Supabase Law Substrate defines the data-layer governance infrastructure
through which all capital-flow actors (donors, ministry administrators, event
participants) interact with the system. The RLS Authorization Court establishes
that the database, not the frontend, is the authoritative host of capital-flow
data. The five demonstration-mode tables (actor\_commands, actor\_events,
actor\_receipts, actor\_outbox, actor\_quarantine) are the specific technical
surfaces through which capital-flow actors emit and consume process events. Their
current ``Allow public access'' policy is a flagged liability; their production
hardening is a receipted next step.

The ZoeApp Ministry Ontology encodes the capital-flow domain itself: Connect
(membership and relationship), Give (financial transaction and stewardship),
Watch (content consumption and engagement), Serve (volunteer and operational
participation), and Pray (liturgical and intercessory participation) are the
five capital-flow channels of a ministry organization. By encoding these as
process-laws rather than UI categories, the atlas establishes that each channel
can be receipted, replayed, audited, and projected onto a board-admissible
claim. This is the contribution to capital-flow infrastructure: the ministry
is a receipted process, not an observed phenomenon.

The Reverse Porter Five Algebra (from the SPR thesis) applies directly here:
as competitive pressure on ministry organizations increases (substitutes,
new entrants, donor bargaining power), the demand for validation authority
--- the ability to prove that ministry capital flowed through lawful, receipted
channels --- increases correspondingly.
